\subsection{BAIRE SPACES}

DEFINITION 3. A topological space X is said to be a Baire space if it satisfies one or the other of the following two equivalent conditions:

(EB) Every countable intersection of dense open sets in \(\mathbf{X}\) is dense in \(\mathbf{X}\) .

(EB') Every countable union of closed sets with no interior points in X has no interior point in X.

Axiom (EB) can be stated in two other equivalent forms:

(EB'') Every non-empty open set in X is non-meagre.

Indeed a set is meagre if and only if it is contained in a countable union of closed sets with no interior points.

(EB''\,``) The complement of a meagre set in X is dense in X.

This signifies that a meagre set cannot contain a non-empty open set, and is therefore equivalent to \((EB'')\) .

PROPOSITION 3. Every non-empty open subspace Y of a Baire space X is a Baire space.

This follows from (EB''), for every open (resp. meagre) set in Y is open (resp. meagre) in X.

It follows from Proposition 3 that every point of a Baire space has a fundamental system of neighbourhoods, each of which is a Baire space. Conversely:

PROPOSITION 4. If every point of a topological space X has a neighbourhood which is a Baire space, then X is a Baire space.

Let A be a non-empty open set in X, let \(x \subset A\) and let V be an open neighbourhood of x which is a Baire space. If A were meagre in X, then \(V \cap A\) would be meagre in V and open in V, which is contrary to hypothesis.

PROPOSITION 5. In a Baire space X, the complement of a meagre set is a Baire space.

Let A be a meagre set in X; then its complement Y = CA in X is dense in X. Let B be a meagre set relative to Y; B is also meagre relative to X, hence A ∪ B is meagre relative to X. Hence the complement of A ∪ B in X, which is also the complement of B in Y, is dense in X and a fortiori dense in Y. Hence Y is a Baire space.

THEOREM I (Baire). (i) Every locally compact space X is a Baire space. (ii) Every topological space X on which there exists a metric, compatible with the topology of X and which defines on X the structure of a complete metric space, is a Baire space.

We shall show that axiom (EB) is satisfied in each case. Let \((\mathbf{A}_n)\) be a sequence of dense open sets in \(\mathbf{X}\) , and let \(\mathbf{G}\) be any non-empty open set. We can then define inductively a sequence \((\mathbf{G}_n)\) of non-empty open sets such that \(\mathbf{G}_1 = \mathbf{G}\) and \(\overline{\mathbf{G}}_{n+1} \subset \mathbf{G}_n \cap \mathbf{A}_n\) : for since by hypothesis \(\mathbf{G}_n\) is not empty, \(\mathbf{G}_n \cap \mathbf{A}_n\) is a non-empty open set; and since \(\mathbf{X}\) is regular in both cases envisaged, there is a non-empty open set \(\mathbf{G}_{n+1}\) such that \(\overline{\mathbf{G}}_{n+1} \subset \mathbf{G}_n \cap \mathbf{A}_n\) . Hence the set \(\mathbf{G} \cap \bigcap_{n=1}^{\infty} \mathbf{A}_n\) contains the intersection of the sets \(G_{n}\) , which is equal to the intersection of the sets \(\overline{G}_{n}\) ; hence it is enough to show that \(\bigcap_{n=1}^{\infty}\overline{G}_{n}\neq\varnothing\) . Now if X is locally compact, we may assume that \(\overline{G}_{2}\) is compact; in the compact space \(\overline{G}_{2}\) , the \(\overline{G}_{n}(n\geqslant2)\) form a decreasing sequence of non-empty closed sets, and their intersection is therefore not empty by axiom (C'\,') (cf.~Chapter I, §9, no. 1, Definition 1{]}. If X is a complete metric space (with respect to a metric compatible with its topology) we may suppose that \(\overline{G}_{n}\) has been chosen so that its diameter (with respect to this metric) tends to o as n tends to \(+\infty\) ; the \(\overline{G}_{n}\) therefore form a Cauchy filter base which converges to a point x, and x necessarily belongs to the intersection of the sets \(\overline{G}_{n}\) . Q.E.D.

Remark. There are Baire spaces which belong to neither of these two categories, in particular Baire spaces which are neither metrizable nor locally compact (Exercise 16); there are also metrizable Baire spaces which possess no complete metric space structure compatible with their topology (Exercise 14).

